\ExerciseSection

\begin{enumerate}
\def\labelenumi{\arabic{enumi})}
\item
  试证明 \(\mathbf{R}^n\) 与 \(\mathbf{R}\) 等势，利用康托尔定理（第 IV 章，§ 8，第 6 小节，定理 1）以及关系式 \(2^{\mathfrak{a}} \cdot 2^{\mathfrak{a}} = 2^{\mathfrak{a}}\) （它对一切无限基数 \(\mathfrak{a}\) 成立）。
\item
  存在区间 \(\mathbf{I} = [\mathrm{o},\mathrm{i}]\) （ \(\mathbf{R}\) 的）到正方形 \(\mathbf{I}\times \mathbf{I}\) （ \(\mathbf{R}^2\) 的）上的连续映射（“佩亚诺曲线”，参见习题 8）。（首先借助第 IV 章，§ 8 的习题 11 证明存在连续映射 \(f\) ，它把康托尔三分集 \(\mathbf{K}\) 映到 \(\mathbf{I}\times \mathbf{I}\) 上，然后把 \(f\) 延拓到 \(\mathbf{I}\) 上。）
\item
  设 A 与 B 是 \(R^{2}\) 的两个可数稠密子集。证明：存在 \(R^{2}\) 到其自身的同胚，它把 A 变换成 B。（首先证明：借助一个旋转，可以化归为投影 \(pr_{1}\) 与 \(pr_{2}\) 都是 A 与 B 到 R 中的单射映射的情形。然后，用一个适当的归纳程序，定义 A 到 B 上的一个双射，它决定 \(pr_{1}A\) 到 \(pr_{1}B\) 上的一个单调映射以及 \(pr_{2}A\) 到 \(pr_{2}B\) 上的一个单调映射。最后，利用第 IV 章，§ 2 的习题 11，证明这个映射是一个同胚，并且可以延拓为 \(R^{2}\) 到其自身的同胚。）由此推出：若 C 是 \(R^{2}\) 的一个子集，它的补集稠密，则 C 同胚于 \(Q^{2}\) 在 \(R^{2}\) 中的补集的一个子集。把这些结果推广到 \(R^{n}\) ，n \textgreater{} 2。
\item
  \(\mathbf{R}^n\) 的每个没有孤立点的可数子空间都同胚于有理直线 \(\mathbf{Q}\) （应用第 IV 章，§ 8 的习题 13）。
\item
  \(\mathbf{R}^n\) 的每个没有孤立点的全不连通紧子空间都同胚于康托尔三分集（利用第 IV 章，§ 8 的习题 11 与习题 12，以及第 II 章，§ 4，第 4 小节的命题 6）。
\item
  \(R^{n}\) 的子集 L 称为折线，如果存在点的有限序列 \((x_{i})_{0\leqslant i\leqslant p}\) （由 \(R^{n}\) 中的点组成），使得：若用 \(S_{i}\) 表示以 \(x_{i-1}\) 与 \(x_{i}\) 为端点的线段（ \(i\leqslant i\leqslant p\) ），则 L 是诸 \(S_{i}\) 的并，这些线段称为 L 的边。（一般地，定义同一条折线的 \(R^{n}\) 中点的有限序列有无穷多个。）折线也是 {[}o, i{]} 到 \(R^{n}\) 中的映射 u 的象：存在一个严格递增序列 \((t_{j})_{0\leqslant j\leqslant q}\) （在 {[}o, i{]} 中取值），\(t_{0}=0\) 且 \(t_{q}=1\) ，使得 u 是仿射线性映射
\end{enumerate}

\[
t \rightarrow \boldsymbol {a} _ {j} + t \boldsymbol {b} _ {j} \quad \text { for } \quad t _ {j - 1} \leqslant t \leqslant t _ {j} \quad \text { and } \quad i \leqslant j \leqslant q
\]

（这样的映射 \(u\) 称为“分段线性的”。）给定 \(\mathbf{R}^n\) 的一个非空子集 A，我们说 A 的两点 \(a, b\) 在 A 中可用折线连结，如果存在折线 \(L \subset A\) ，它由一个序列 \((x_i)_{0 \leqslant i \leqslant p}\) 定义，使得 \(x_0 = a\) 且 \(x_p = b\) 。若 A 的任意两点都能在 A 中用折线连结，则 A 是连通的。反之，若 A 是 \(\mathbf{R}^n\) 的连通开子集，证明 A 的任意两点都能在 A 中用折线连结（考察关系“ \(x\) 与 \(y\) 在 A 中可用折线连结”（它是 A 的点 \(x, y\) 之间的关系）；证明这是一个等价关系，而且它的类都是开集）；甚至可以始终假设这条折线是一些线段的并，其中每条线段都平行于某条坐标轴（同样的方法）。由此推出：若 A 是 \(\mathbf{R}^n\) 中的非空开集，则点 \(a \in A\) 的分支是 A 中能与 \(a\) 用 A 中折线连结的一切点组成的集。

\begin{enumerate}
\def\labelenumi{\arabic{enumi})}
\setcounter{enumi}{6}
\item
  设 A 是 \(R^{n}\) 中的非空连通开集（n \textgreater{} 1），并设 \((\mathrm{V}_{p})_{p \in \mathbb{N}}\) 是 \(R^{n}\) 中线性簇的一个可数族，其中每个线性簇的维数都 \(\leqslant n - 2\) 。证明：若 B 是诸 \(V_{p}\) 的并，则 \(A \cap \mathfrak{C}B\) 在 A 中稠密且连通。（为证明 \(A \cap \mathfrak{C}B\) 在 A 中稠密，对 n 用归纳法。然后证明：若 x, y 是 A 的两个不同的点，则对每个 \(\varepsilon > 0\) 存在点 \(y' \in A\) 满足 \(||y' - y|| \leqslant \varepsilon\) ，并且以 x 与 \(y'\) 为端点的闭线段与 B 只交于点 x。最后利用习题 6 证明 \(A \cap \mathfrak{C}B\) 连通。）
\item
  由习题 7 推出：若 \(n > \mathbf{I}\) ，则 \(\mathbf{R}^n\) 的非空开子集不与 \(\mathbf{R}\) 的任何子集同胚（*）。
\item
  \(\mathbf{R}^{n}\ (n > \mathrm{i})\) 中的折线 L（习题 6）称为简单折线，如果它同胚于 R 的区间 \(I = [0, i]\)。这等价于说存在 I 到 L 上的一个分段线性同胚 u（习题 6）。证明：若 A 是 \(\mathbf{R}^2\) 中的连通开集，且 \(\mathbf{L}\subset\mathbf{A}\) 是简单折线，则 \(\mathbf{A}-\mathbf{L}\) 连通（对组成 L 的线段的个数用归纳法论证，并利用第 I 章，§ 11 的习题 6）。
\end{enumerate}

* 10) a) 设 M 是有限集，R \(x, y\) 是 M 的元素 x, y 之间的一个对称关系。假设 M 中存在两个不同的元素 a, b，它们具有下列性质：存在 x ≠ a 与 y ≠ b，使得 x（相应地 y）是使 R \(a, z\) （相应地 R \(b, z\) ）为真的唯一的元素 z ∈ M；并且对 M 中每个异于 a 与 b 的 t，使 R \(t, z\) 为真的所有 z ∈ M 组成的集是一个由两个都异于 t 的元素组成的集。证明：在这些条件下，存在一个双射 i → x\_i （从 N 的区间 {[}o, n{]} 到 M 上），使得 x\_0 = a，x\_n = b，并且对 i ≤ i ≤ n，R \(x_{i-1}, x_i\) 为真（对 i 用归纳法定义 x\_i）。

\begin{enumerate}
\def\labelenumi{\alph{enumi})}
\setcounter{enumi}{1}
\tightlist
\item
  在 \(\mathbf{R}^n\) （ \(n \geqslant 2\) ）中——把它与 \(\mathbf{R}^{n-1} \times \mathbf{R}\) 等同起来——设 \(\mathbf{L}\) 与 \(\mathbf{L}'\) 是两条折线，分别由序列 \((\boldsymbol{x}_i)_{0 \leqslant i \leqslant p}\) 与 \((\boldsymbol{x}_j')_{0 \leqslant j \leqslant q}\) 定义，且具有下列性质：(i) 若记 \(\boldsymbol{x}_i = (\boldsymbol{y}_i, z_i)\) ，\(\boldsymbol{x}_j' = (\boldsymbol{y}_j', z_j')\) ，则我们有 \(z_0 = z_0' = 0\) ，\(z_p = z_q' = 1\) ，且 \(0 < z_i < 1\) 对 \(1 \leqslant i \leqslant p - 1\) 成立，\(0 < z_j' < 1\) 对 \(1 \leqslant j \leqslant q - 1\) 成立；(ii) \(p + q - 2\) 个数 \(z_i, z_j'\) （ \(1 \leqslant i \leqslant p - 1\) ，\(1 \leqslant j \leqslant q - 1\) ）两两不同；(iii) \(\mathbf{L}\) （相应地 \(\mathbf{L}'\) ）的两条不同的边至多有一个公共点{[}这个点可以是、也可以不是诸 \(\boldsymbol{x}_i\) （相应地 \(\boldsymbol{x}_j'\) ）之一{]}。证明：存在两个满射的分段线性映射（习题 6）\(\boldsymbol{u}\) ： \(\mathbf{I} \to \mathbf{L}\) ，\(\boldsymbol{u}'\colon \mathbf{I} \to \mathbf{L}'\) ，其中 \(\mathbf{I}\) 是区间 \([0, 1]\) （ \(\mathbf{R}\) 中的），使得：若记 \(\boldsymbol{u}(t) = (\boldsymbol{v}(t), \zeta(t))\) ，\(\boldsymbol{u}'(t) = (\boldsymbol{v}'(t), \zeta'(t))\) （ \(t \in \mathbf{I}\) ），则我们有 \(\zeta(t) = \zeta'(t)\) （对一切 \(t \in \mathbf{I}\) ），
\end{enumerate}

\[
\zeta (\mathrm{o}) = \zeta^ {\prime} (\mathrm{o}) = \mathrm{o} \quad \text { and } \quad \zeta (\mathrm{i}) = \zeta^ {\prime} (\mathrm{i}) = \mathrm{i}.
\]

{[}设 \((a_{k})_{1\leqslant k\leqslant p + q - 2}\) 是由异于 o 与 i 的诸 \(z_{i}\) 与诸 \(z_j^{\prime}\) 组成的严格递增序列，并令 \(a_0 = 0\) ，\(a_{p + q - 1} = \mathrm{i}\) ；设 \(\mathbf{B}_i\) 表示由满足下列条件的所有 \(\pmb {x} = (\pmb {y},z)\in \mathbb{R}^n\) 组成的集：

\[
a _ {i - 1} \leqslant z \leqslant a _ {i} \quad (\mathrm{I} \leqslant i \leqslant p + q - \mathrm{I});
\]

设 \(\mathfrak{M}\) 表示所有偶 \(\gamma = (\mathrm{C},\mathrm{C}^{\prime})\) 组成的集，其中 \(\mathrm{C}\) （相应地 \(\mathrm{C}'\) ）是同一个 \(\mathbf{B}_i\) 与 \(\mathbf{L}\) （相应地 \(\mathbf{L}'\) ）的一条边的交，且 \(\mathbf{C}\) 与 \(\mathbf{C}'\) 都不由单独一个点组成；又设 \(\mathbf{R}\wr \alpha, \beta\wr\) 表示元素 \(\alpha, \beta\) （ \(\mathfrak{M}\) 的）之间的下述关系：“ \(\alpha \neq \beta\) ；若 \(\alpha = (\mathrm{C}_1,\mathrm{C}_1')\) ，\(\beta = (\mathrm{C}_2,\mathrm{C}_2')\) ，则 \(\mathrm{C}_1 \cap \mathrm{C}_2\) 与 \(\mathrm{C}_1' \cap \mathrm{C}_2'\) 都非空，其中之一由单独一个点组成，而且若这个点不是诸 \(x_i\) （相应地 \(x_j'\) ）之一，则 \(\mathrm{C}_1\) 与 \(\mathrm{C}_2\) （相应地 \(\mathrm{C}_1'\) 与 \(\mathrm{C}_2'\) ）都含于 \(\mathbf{L}\) （相应地 \(\mathbf{L}'\) ）的同一条边中。证明 \(a\) ) 可应用于关系 \(\mathbf{R}.\) {]}

\begin{enumerate}
\def\labelenumi{\alph{enumi})}
\setcounter{enumi}{2}
\tightlist
\item
  在 \(\mathbf{R}^n\) 中（把它与 \(\mathbf{R}^{n-1} \times \mathbf{R}\) 等同起来），设 \(\mathbf{K}\) 是连通紧集，使得 \(\mathbf{K} \cap (\mathbf{R}^{n-1} \times \{\mathbf{o}\})\) 至少含两个不同的点。证明：若 \(\mathbf{K}' = \mathbf{K} - \mathbf{K}\) （即所有 \(x - y\) 组成的集，其中 \(x \in \mathbf{K}\) 且 \(y \in \mathbf{K}\) ），则 \(\mathbf{K}' \cap (\mathbf{R}^{n-1} \times \{\mathbf{o}\})\) 含有一个不由单独一个点组成的连通集。{[}利用 b)，并应用第 II 章，§ 4 的命题 6 与习题 15。{]}
\end{enumerate}

\begin{enumerate}
\def\labelenumi{\arabic{enumi})}
\setcounter{enumi}{10}
\tightlist
\item
  把第 IV 章，§ 8 的习题 14 与习题 15 推广到空间 \(\mathbf{R}^n\) （ \(n \geqslant 2\) ）。
\end{enumerate}

* 12) a) 设 \(f\) 是 \(\mathbf{R}\) 到 \(\mathbf{R}\) 中的映射，并设 \(\mathbf{G} \subset \mathbf{R}^2\) 是 \(f\) 的图象。假设 \(\mathbf{G}\) 在 \(\mathbf{R}^2\) 中稠密，并且与 \(\mathbf{R}^2\) 中每个不完全含于形如 \(\{x_n\} \times \mathbf{R}\) 的直线的可数并中的完全集相交。证明 \(\mathbf{G}\) 连通。（若不然，证明：存在两个非空的不相交开集 \(\mathbf{A}\) ，\(\mathbf{B}\) （ \(\mathbf{R}^2\) 中的），使得 \(\mathbf{G}\) 是 \(\mathbf{G} \cap \mathbf{A}\) 与 \(\mathbf{G} \cap \mathbf{B}\) 的并，\(\mathbf{A}\) 中至少有一点与 \(\mathbf{B}\) 中的一点有相同的横坐标（利用 \(\mathbf{R}\) 连通这一事实）；然后利用习题 11。）

\begin{enumerate}
\def\labelenumi{\alph{enumi})}
\setcounter{enumi}{1}
\tightlist
\item
  由 a) 推出：存在 R 到 R 中的线性映射 f，它的图象 G 在 \(R^{2}\) 中稠密且连通。{[}利用 a) 与习题 11，按照《集合论》第 III 章，§ 6 习题 24 中的同样方法，用超限归纳法定义 f 在 R 的一个哈梅尔基的各点上的值。{]}由此推出：\(R^{2}\) 的子群 G 满足第 V 章，§ 3 习题 4 中的条件 ( \(R_{I}\) ) 与 ( \(R_{II}\) ) ，但不是局部紧的，因而它的拓扑严格细于 R 的通常拓扑；并且 G 不是局部连通的。
\end{enumerate}

\begin{enumerate}
\def\labelenumi{\arabic{enumi})}
\setcounter{enumi}{12}
\tightlist
\item
  把 \(R^{n}\) 与 \(R^{n-1} \times R\) 等同起来，设 f 是 \(R^{n-1}\) 的一个开子集 A 到 R 中的连续映射，并设 S 是它的图象。证明：\(R^{n}\) 的子空间 S 同胚于 A，并且 S 在“圆柱” \(A \times R\) 中的补集是 \(A \times R\) 中的稠密开集，而且当 A 连通时它恰有两个分支。
\end{enumerate}

